\subsection{The link comparison in Lead 7}

\begin{proposition}[Proved: the link comparison, with its necessary hypothesis]\label{C:link-comparison}
Let $A$ and $D$ be small categories, let $u:D\to A$ be a functor, and let
\[
 R=\mathop{\rm colim}_{d\in D}A(-,u(d)):A^{\rm op}\longrightarrow\mathrm{Set}.
\]
Assume that every connected component of every comma category $x\downarrow u$
is weakly contractible. Form $A'$ by adjoining one object $Z$, by setting
$A'(x,Z)=R(x)$ for $x\in A$, by allowing no arrows from $Z$ to $A$, and by
setting $A'(Z,Z)=\{1_Z\}$; composition with arrows of $A$ is the presheaf
action of $R$. Let $H':A'\to\mathrm{Set}$ be any covariant functor, and put
$H=H'|_A$.

Define the two-sided category $E(R,H)$ to have objects $(x,r,k)$ with
$x\in A$, $r\in R(x)$ and $k\in H(x)$. An arrow
$(x,r,k)\to(y,s,l)$ is an arrow $b:x\to y$ of $A$ satisfying
$R(b)s=r$ and $H(b)k=l$. Then the functor
\[
 F:\operatorname{El}_D(u^*H)\longrightarrow E(R,H),\qquad
 (d,k)\longmapsto(u(d),[d,1_{u(d)}],k)
\]
which sends an arrow $\delta:d\to d'$ to $u(\delta)$ is a weak equivalence
on nerves. For every $h\in H'(Z)$ the old comma link
\[
 L_h=(\operatorname{El}_A H)\downarrow (Z,h)
\]
is a full union of connected components of $E(R,H)$. If
$e=\dim ND<\infty$, then $H_j(BL_h;\mathbb Z)=0$ for $j>e$ and the
inclusion induces an injection
\[
 H_j(B\operatorname{El}_A H;\mathbb Z)\longrightarrow
 H_j(B\operatorname{El}_{A'}H';\mathbb Z)
 \quad(j>e).
\]
The assertion also holds with arbitrary constant coefficients. If $D$ is
empty, $R$ and all links are empty and adjoining $Z$ adds only isolated
objects to the category of elements.
\end{proposition}

\begin{proof}
For any $x\in A$, the set $R(x)$ is canonically
$\pi_0(x\downarrow u)$: its representatives are pairs $(d,b:x\to u(d))$,
and its generating identifications are
$(d,b)\sim(d',u(\delta)b)$ for $\delta:d\to d'$.
Fix $w=(x,r,k)\in E(R,H)$. An object of $w\downarrow F$ is specified by
$d\in D$, $l\in H(u(d))$ and $b:x\to u(d)$ with
$[d,b]=r$ and $H(b)k=l$. Thus $l$ is forced. Sending this object to
$(d,b)$ gives an isomorphism of categories
\[
 w\downarrow F\ \cong\ (x\downarrow u)_r,
\]
where the right side is the connected component labelled by $r$. Indeed,
an arrow on either side is exactly a $\delta:d\to d'$ with
$u(\delta)b=b'$. The assumed weak contractibility proves the nerve comparison as follows.
For a general functor $F:X\to Y$ with these comma nerves weakly
contractible, form the bisimplicial set
\[
 B_{p,q}=\{(y_0\to\cdots\to y_p\to F(x_0),\
                 x_0\to\cdots\to x_q)\}.
\]
For a fixed $q$-simplex of $NX$, its horizontal fiber is
$N(Y\downarrow F(x_0))$, a nerve with terminal object
$(F(x_0),1)$. For a fixed $p$-simplex of $NY$, its vertical fiber is
$N(y_p\downarrow F)$, weakly contractible by hypothesis. Realizing first
in either direction therefore shows that the projections from the
diagonal realization to $BX$ and $BY$ are weak equivalences. Here the
levelwise statement passes to realization because degeneracy inclusions
of simplicial sets are cellular inclusions; the usual skeletal
filtration proves the assertion by successive gluing along those
inclusions. The two maps to $BY$, one by its projection and one by
$BF$ after projection to $BX$, are homotopic: the joining arrow
$y_p\to F(x_0)$ is the joining edge of the displayed chain and supplies
the simplicial mapping-cylinder homotopy. Thus $BF$ is a weak
 equivalence. Applying this argument to the comma categories just
computed proves the claimed comparison.

The function
\[
 \lambda:\operatorname{Ob}E(R,H)\longrightarrow H'(Z),\qquad
 (x,r,k)\longmapsto H'(r)k
\]
is constant on every arrow, by functoriality. The full subcategory
$\lambda^{-1}(h)$ is exactly $L_h$, including its arrows. It is consequently
a union of connected components. The restriction of $F$ to the union of
source components mapping into $L_h$ remains a weak equivalence: the
same comma-category computation applies separately on these components.
The projection $\operatorname{El}_D(u^*H)\to D$ is a discrete opfibration.
A nonidentity arrow cannot project to an identity, so a nondegenerate
simplex in its nerve has dimension at most $e$. This gives the homology
vanishing for every link.

For completeness, the attachment assertion is a chain-level statement,
not an assumption that a link embeds as a subcategory of the old category.
Every nondegenerate simplex of $N\operatorname{El}_{A'}H'$ that is not
old ends at exactly one $(Z,h)$; deleting this final vertex gives a
simplex of the corresponding comma link, including the final arrow.
Hence the normalized relative chain complex is canonically
\[
 C_*(N\operatorname{El}_{A'}H',N\operatorname{El}_A H)
 \cong\bigoplus_{h\in H'(Z)}\widetilde C_{*-1}(NL_h).
\]
Its long exact sequence contains
\[
 \bigoplus_h H_j(BL_h)\longrightarrow H_j(B\operatorname{El}_A H)
 \longrightarrow H_j(B\operatorname{El}_{A'}H'),
\]
and the first group is zero for $j>e$. All chains have finite support,
so the direct sum and the argument remain valid for infinite sets of
new elements and for arbitrary small indexing categories.
\end{proof}

\begin{proposition}[Proved: the equalizing-precomposition obstruction]\label{C:link-obstruction}
Consider a cyclic quasi-filter indexing graph with copies of $S$ at its
vertices and copies of $T$ on its edges, the two maps on each edge being
chosen from a parallel pair $f,g:T\rightrightarrows S$. If an arrow
$c:x\to T$ satisfies $fc=gc$, the component of $x\downarrow u$ containing
these lifts has nonzero first integral homology. If $f$ and $g$ are
monomorphisms, this component is precisely a subdivided circle.
Consequently the comma hypothesis in Proposition~\ref{C:link-comparison}
fails. In the repaired first-stage AKR quasi-filter for $v_1,v_2$, the
last and first maps satisfy $\alpha_{4k}v_1=\alpha_0v_1$, so this
obstruction occurs when quasi-filtering $(\alpha_0,\alpha_{4k})$ at the
next stage.
\end{proposition}
\begin{proof}
At every vertex choose the arrow $fc=gc:x\to S$, and at every edge
choose $c:x\to T$. These choices form a copy of the subdivided indexing
circle in $x\downarrow u$. This comma category has no nonidentity
composable pair, because its projection to the bipartite indexing graph
has none. Its nerve is therefore a graph; the displayed cycle cannot
be the boundary of a two-chain and has nonzero first homology. If $f$
and $g$ are monomorphisms, either map $A(x,T)\to A(x,S)$ is injective.
The vertices of the displayed cycle then have no additional incident
edges, and its component is exactly that circle. The final application
is the last neighboring equality in the repaired formula, with
$c=v_1$. In a category where $f,g$ are not monomorphisms the component
can have further branches or cycles; only the noncontractibility, not
literal equality with a circle, is automatic.
\end{proof}

\paragraph{What this does not prove about the infinite AKR category.}
The dimension-one conclusion applies only when the strict colimit of
representables has the stated comma-category comparison. The preceding
proposition supplies explicit AKR attachments where that hypothesis
fails. Thus it cannot, by itself, establish persistence of the
first-stage two-dimensional class. In particular we do not infer
$H_2(L_h)=0$ for those exceptional links from the dimension of the
quasi-filter indexing graph.

\begin{proposition}[Proved: all-ring consequence of the known counting invariant]\label{C:akr-all-rings}
For every admissible choice in the repaired AKR Construction C.1, for
every nonzero unital commutative ring $R$, and for every object $X$ of
$K^*$, the $R$-linear map
\[
 R[K^*(W,X)]\longrightarrow R[K^*(V,X)]^2,\qquad
 t\longmapsto\bigl(t(v_2-v_1),t(v_3-v_1)\bigr)
\]
is injective. In particular, the three initial point maps have no
$R$-affine equalizer. The same injectivity holds for finitely supported
coefficients in an arbitrary abelian group, with the evident meaning
of the two differences.
\end{proposition}
\begin{proof}
We use the counting invariant of the opened AKR Theorem C.2,
pp.~385--389, with the repaired maps satisfying its required (9)--(11).
We checked the induction, not merely its conclusion. Its statements are:
for an object $X$ created by stage $m$ and an arrow $g:Y\to X$,
\begin{enumerate}
\item if a $W$-arrow $h:W\to X$ has at least two image points in the
image of $g$, it factors through $g$;
\item if $h_1,\ldots,h_r:W\to X$ are distinct and each meets the image
of $g$ in at most one point, then their union outside that image has
at least $a_mr$ points.
\end{enumerate}
Here $3/2<a_m\leq2$. To check the source's inductive use of these
statements, write every incoming map at a new quasi-filter target as
$\alpha_i$ followed by an old map. Property (10) moves any overlap of
two copies through consecutive copies, and (11) identifies each
consecutive overlap with the image of one of the old pair. The first
inductive statement then routes the entire $W$-arrow through the
neighbor, proving the new first statement. For the second statement,
partition the $W$-arrows by their first available copy, discard a block
of $3n$ consecutive indices of least cardinality, and sweep the
remaining indices from the two ends toward that block. Every newly
counted $W$-arrow meets the preceding overlap in at most one point:
otherwise the just-proved routing statement would have assigned it to
an earlier copy. The old second statement thus adds at least $a_m$
points per retained arrow. The discarded fraction is at most
$3n/s\leq1-a_{m+1}/a_m$, giving $a_{m+1}r$. A disjoint-sum target
separates the two sets of arrows; the same inequalities apply separately
to its summands. These are precisely the induction steps on
pp.~386--389, and all use of overlaps is justified by the repaired
(10)--(11), not the invalid restriction assertion.

For clarity, the source's displayed set $L=\{1,\ldots,n\}\setminus
\widetilde H_{j_0}$ on p.~387 must read $\{1,\ldots,r\}$; all surrounding
cardinality formulas use the latter set. This is an index typo, not an
additional hypothesis.

Fix now a nonempty finite set $S$ of $r$ distinct $W$-arrows into $X$.
Two distinct such arrows meet in at most one point: otherwise the first
statement would factor one through the other via an endomorphism of
$W$, and $K^*(W,W)=\{1_W\}$. Applying the second statement with one
member of $S$ as $g$ therefore gives
\[
 \left|\bigcup_{h\in S}h(W)\right|
 \geq a_m(r-1)+3>\frac32r.
\]
This is also valid for $r=1$. There are exactly $3r$ incidences
$(h,i)$ with $h\in S$ and $i\in\{1,2,3\}$. If each image point had
at least two incidences, the displayed union would have at most
$3r/2$ points. Hence some image point has exactly one incidence.

Let $t=\sum_{h\in S}c_hh$ have all $c_h\ne0$, and choose that private
point, say $h_0v_i$. Its coefficient in $tv_i$ is $c_{h_0}$, whereas
its coefficient in $tv_j$ is zero for either $j\ne i$. Thus the three
vectors $tv_1,tv_2,tv_3$ cannot be equal. This uses no division and no
restriction on the characteristic or on zero divisors. It proves the
claimed injectivity for coefficients in any abelian group and hence
for every specified ring. In particular, it applies to
$\mathbb Q$ and every $\mathbb F_p$ simultaneously.
\end{proof}

\begin{proposition}[Proved: an exact cycle-space test for one attachment]\label{C:attachment-cycle-test}
Let $A'$ be obtained from $A$ by adjoining one new maximal object $Z$
with identity endomorphism only. Suppose a family
$\alpha_i:T_i\to Z$ ($i\in I$, possibly infinite) generates all incoming
arrows: every $b:X\to Z$, $X\in A$, can be written
$\alpha_i c$ with $c:X\to T_i$ in $A$. For any covariant functor
$F':A'\to\mathrm{Ab}$, write $F=F'|_A$ and
$Q=\operatorname{colim}_A F$. Define homomorphisms
\[
 a:\bigoplus_{i\in I}F(T_i)\longrightarrow F'(Z),\quad
   (x_i)\longmapsto\sum_i F'(\alpha_i)x_i,
 \qquad
 b:\bigoplus_{i\in I}F(T_i)\longrightarrow Q,\quad
   (x_i)\longmapsto\sum_i[x_i].
\]
Then, canonically,
\[
 \operatorname{colim}_{A'}F'
   \cong (Q\oplus F'(Z))/\{(b(x),-a(x)):x\in\bigoplus_iF(T_i)\},
\]
\[
 \ker(Q\to\operatorname{colim}_{A'}F')=b(\ker a).
\]
For the AKR action
$H=h_V/(cv_1\sim cv_2)$, put
$K(X)=\ker(\mathbb Z[A(W,X)]\xrightarrow{,v_2-v_1,}
\mathbb Z[A(V,X)])$. Thus $K(X)$ is the integral cycle group of the
graph whose vertices are $V$-arrows and whose edges are $W$-arrows,
with edge $c$ directed from $cv_1$ to $cv_2$. For a quasi-filter
attachment all $T_i=T$. Its old-to-new map on $H_2(B\operatorname{El}H)$
is injective if and only if
\[
 \sum_i\alpha_i\kappa_i=0\ \text{in }K(Z)
 \quad\Longrightarrow\quad
 \left[\sum_i\kappa_i\right]=0\ \text{in }\operatorname{colim}_A K
\]
for every finitely supported family $\kappa_i\in K(T)$.
\end{proposition}
\begin{proof}
A cocone from $F'$ to an abelian group $M$ is an old cocone
$Q\to M$ together with a map $F'(Z)\to M$ agreeing along all incoming
arrows. Agreement along $\alpha_i$ suffices: for
$b=\alpha_i c$, naturality of the old cocone gives agreement along
$c$ first, hence along $b$. This proves the displayed pushout
presentation. The class of $(q,0)$ is zero precisely when
$(q,0)=(b(x),-a(x))$ for some $x$, which proves the kernel formula.
No ordinary single-object representative for a general colimit element
is assumed here; the direct sum explicitly allows finite sums at
several objects.

For the last assertion, linearizing the binary equivalence relation
gives
\[
 0\longrightarrow K\longrightarrow P_W\xrightarrow{,v_2-v_1,}
 P_V\longrightarrow\mathbb Z[H]\longrightarrow0.
\]
Split this at the image of the middle map. The projectivity of $P_W$
and $P_V$, together with
$\operatorname{colim}P_W=\operatorname{colim}P_V=\mathbb Z$ and the
zero induced middle map, gives
\[
 H_2(B\operatorname{El}_A H;\mathbb Z)
 \cong\ker(\operatorname{colim}_A K\xrightarrow{\rm aug}\mathbb Z).
\]
This is the binary instance of the formula whose proof is given in
TOTEX Part B; we checked it by these two short exact sequences and by
the identification of the simplicial replacement of $\mathbb Z[H]$
with the nerve chain complex of $\operatorname{El}_A H$. These
identifications are natural under the stated extension, since the old
representables and their relation map are its restrictions. The
augmentation of an element in the kernel of the old-to-new map on
$\operatorname{colim}K$ is necessarily zero. Hence that whole kernel
is precisely the kernel on $H_2$. Applying the first part with
$F'=K$ proves the criterion.
\end{proof}

\begin{example}[Refuted literal wording: without monomorphisms the component need not be a circle]\label{C:nonmonic-circle-counterexample}
The unqualified phrase ``the component is a circle'' in Lead~7 is
false if it is read without the monomorphism hypothesis of the AKR
application. Let $A$ have objects $x,T,S$, arrows
\[
 A(x,T)=\{c,d\},\qquad A(T,S)=\{f,g\},\qquad A(x,S)=\{h\},
\]
identity endomorphisms only, no other arrows, and the four compositions
$fc=fd=gc=gd=h$. Composition with identities is prescribed in the
usual way. This is a complete finite category: there are no three
composable nonidentity arrows, and the identity laws settle every
remaining associativity equation.

Take the five-vertex cyclic quasi-filter indexing graph ($k=1$).
It has copies $S_i$ of $S$ and copies $T_i$ of $T$, with two arrows
from $T_i$ to $S_i,S_{i+1}$ (cyclic indices), mapped by $u$ to
$(g,f)$ for $i=0,1$, to $(g,g)$ for $i=2$, to $(f,g)$ for $i=3$,
and to $(f,f)$ for $i=4$. The comma category $x\downarrow u$ has
five $S$-vertices $(S_i,h)$ and ten $T$-vertices
$(T_i,c),(T_i,d)$, and twenty nonidentity arrows. It is connected.
Its nerve is a graph, with $H_1\cong\mathbb Z^{20-15+1}
=\mathbb Z^6$; each $S$-vertex has valence four. Thus its component
is neither isomorphic nor homotopy equivalent to a circle, although
$fc=gc$ and its noncontractibility is exactly as asserted in the
corrected proposition. The complete composition table, all typed
composable triples, and these graph counts are checked by
\texttt{akr\_nonmonic\_comma\_circle.py} and its supplied output.
\end{example}

\begin{proposition}[Refuted: pointwise free split injections do not suffice for acyclicity]\label{C:split-monic-counterexample}
There exists a small nonempty connected quasi-filtered category $A$
and a covariant functor $Q:A\to\mathrm{Ab}$ such that every $Q(X)$ is
finitely generated free abelian, every $Q(g)$ is split injective with
free abelian cokernel, and nevertheless
\[
 T_1(Q)\cong\mathbb Z,\qquad T_i(Q)=0\quad(i\geq2).
\]
One may take any repaired AKR category $A=K^*$ and
\[
 Q=\operatorname{coker}\bigl(D:P_W\longrightarrow P_V^2\bigr),
 \qquad D(t)=\bigl(t(v_2-v_1),t(v_3-v_1)\bigr).
\]
Thus pointwise freeness and split injectivity do not by themselves
close the cycle-module colimit obstruction.
\end{proposition}
\begin{proof}
Proposition~\ref{C:akr-all-rings} says that $D$ is injective, so
$0\to P_W\xrightarrow D P_V^2\to Q\to0$ is a projective resolution
of length one. Its induced colimit map is the zero map
$\mathbb Z\to\mathbb Z^2$, because each component is a difference
of two augmentation-one representable maps. Therefore
$T_1(Q)=\ker(0:\mathbb Z\to\mathbb Z^2)=\mathbb Z$ and all higher
$T_i(Q)$ vanish.

We verify the asserted pointwise and transition properties integrally.
For an arrow $g:Y\to X$, identify the images of $P_V(Y)$ and
$P_W(Y)$ inside their counterparts at $X$, using monicity. The induced
relative map
\[
 \overline D:
 P_W(X)/gP_W(Y)\longrightarrow
 P_V(X)^2/gP_V(Y)^2
\]
has one column for every $W$-arrow $h:W\to X$ not factoring through
$g$. By AKR C.2$(c_m)$, its three-point image meets $g(Y)$ in at most
one point. For any nonempty collection of these columns,
C.2$(d_m)$ supplies more than $3r/2$ points outside $g(Y)$, where $r$
is the number of columns. Hence one such outside point is private to
one column. At that point the column has coordinate pair
$(-1,-1)$, $(1,0)$ or $(0,1)$, according to its position in the triple.
Choose a unit entry from that pair as pivot. Its row has zero entries
in all other columns. Eliminating this column therefore leaves all
other columns unchanged. Repeating gives a full unit-pivot reduction.
All these matrices are finite because each object's incoming hom-sets
are frozen at its finite stage. Thus $\overline D$ is injective and
its cokernel is free abelian. The same argument without the relative
quotient proves that $Q(X)$ is free abelian.

The exact diagram comparing the two short exact sequences at $Y$
and $X$ now gives, by the kernel--cokernel sequence,
\[
 0\longrightarrow Q(Y)\xrightarrow{Q(g)}Q(X)
 \longrightarrow\operatorname{coker}\overline D\longrightarrow0.
\]
The final group is free, so this sequence splits. This verifies every
claimed transition property while retaining the nonzero derived
colimit calculated above.
\end{proof}

\begin{corollary}[Proved: the same obstruction with free groups and free-factor maps]\label{C:free-group-counterexample}
For the same quasi-filtered category $A=K^*$ there is a covariant
functor $\Gamma:A\to\mathrm{Grp}$ such that every $\Gamma(X)$ is a
finitely generated free group and every transition map identifies its
source with a free factor of its target, but the abelianization
functor $\Gamma_{\rm ab}:A\to\mathrm{Ab}$ has
$T_1(\Gamma_{\rm ab})\cong\mathbb Z$.
\end{corollary}
\begin{proof}
For each $X$, take group generators $x_{v,1},x_{v,2}$ indexed by
$v\in A(V,X)$ and impose, for each $h\in A(W,X)$, the relation
\[
 x_{hv_2,1}\,x_{hv_1,1}^{-1}\,
 x_{hv_3,2}\,x_{hv_1,2}^{-1}=1.
\]
Postcomposition on $v$ and $h$ defines a homomorphism on these
presentations for every arrow of $A$, and makes $\Gamma$ a covariant
functor. Its abelianization is exactly the functor $Q$ of
Proposition~\ref{C:split-monic-counterexample}.

For any nonempty family of remaining relations the private-vertex
argument supplies a vertex occurring in just one of them. One of the
two generators indexed by that vertex occurs exactly once in that
relation and nowhere in the others. At $hv_1$ we may choose the
coordinate-one generator; the coordinate-two occurrence is a different
generator. A Tietze elimination solves for the chosen generator and
removes that sole relation, leaving all other relations unchanged.
Repeating removes every relation and proves that $\Gamma(X)$ is free.

For an arrow $g:Y\to X$, retain the entire old presentation at $Y$.
Apply the relative private-vertex argument only to the new relations,
those indexed by $W$-arrows not factoring through $g$. Each chosen
private vertex lies outside the old image, so the eliminated generator
is new and occurs in no old relation. Eliminating all new relations
leaves the old presentation together with additional free generators.
Thus the natural map identifies $\Gamma(Y)$ with a free factor of
$\Gamma(X)$. No compatible choices of free bases or of retractions
are required. Finally the preceding proposition computes the stated
nonzero derived colimit of its abelianization.
\end{proof}
